\section{Splitting vertices and bounding triangle mass}
\label{sec:splitting}

The local walks give similar probability distributions to the two
neighbors at a corner of a triangle, on average under the ordered
triangle law. We now sample from these distributions jointly.
Equal samples keep the two incident edges of a triangle at the same copy
of a vertex. A condition on the reversed transition probability ensures
that each copy has a small weighted neighborhood.

\subsection{Sampling all rows together}

We use the shared rejection construction of Kleinberg and
Tardos~\cite[Section~3, Lemmas~3.1 and~3.2]{KleinbergTardos2002}.
Angel and Spinka~\cite[Theorem~2 and Section~2.1]{AngelSpinka2021}
give the sharper pairwise disagreement estimate stated below.
The short proof supplies the joint law for all rows at once.

\begin{lemma}[Shared rejection sampling]\label{lem:shared-rejection}
Let $I$ and $J$ be finite nonempty sets, and let $(p_i)_{i\in I}$ be
probability distributions on $J$. There are jointly distributed random
variables $(Y_i)_{i\in I}$ such that $Y_i$ has law $p_i$ and, for every
$i,i'\in I$,
\[
 \mathbb P(Y_i\ne Y_{i'})
 \le \frac{2\|p_i-p_{i'}\|_{\mathrm{TV}}}
 {1+\|p_i-p_{i'}\|_{\mathrm{TV}}}
 \le 2\|p_i-p_{i'}\|_{\mathrm{TV}}.
\]
The resulting joint law has finite support.
\end{lemma}

\begin{proof}
Write $m=|J|$. Take independent trials $(Z_t,U_t)_{t\ge1}$, with $Z_t$
uniform on $J$ and $U_t$ uniform on $[0,1]$. For each $i$, let $Y_i=Z_t$
at the first trial satisfying $U_t\le p_i(Z_t)$. A trial is accepted
by row $i$ with probability $1/m$, and its accepted value has law $p_i$.
More explicitly, for $y\in J$,
\[
 \mathbb P(Y_i=y)
 =\sum_{t\ge1}(1-1/m)^{t-1}\frac{p_i(y)}m=p_i(y).
\]
Every row stops almost surely; since $I$ is finite, all rows stop after
a finite number of trials almost surely.

Fix two rows $p,q$, and put $d=\|p-q\|_{\mathrm{TV}}$. At the first
trial accepted by either row, the probability that both accept is
\[
 \frac{\sum_{y\in J}\min(p(y),q(y))}
 {\sum_{y\in J}\max(p(y),q(y))}
 =\frac{1-d}{1+d}.
\]
On this event their labels agree. They may also agree otherwise, so
the disagreement probability is at most $2d/(1+d)$. Finally, the label
vector takes values in the finite set $J^I$, so its induced law has
finite support. The infinite sequence of trials is only a convenient
construction of that finite joint law.
\end{proof}

\subsection{Admissible labels and vertex splitting}

Let $H$ be a finite graph with positive vertex weights. Assume that every
edge $uv$ has common-neighbor weight $c_{uv}\ge1/x$, that every
neighborhood has weight at most $e^x$, and that $H$ satisfies the
cross-mass condition~\eqref{eq:cross-mass} with constant $C\ge1$. Here $x\ge32$. If $H$ has
an edge, then it has a triangle; throughout this subsection suppose
$T_H(w)>0$.

Use the integer $j$ supplied by Lemma~\ref{lem:entropy-smoothing}, and
for each non-isolated vertex $u$ write $K_u=K_u^j$ for the corresponding
kernel on $N_H(u)$. Isolated vertices may be discarded, since they
contribute neither edge nor triangle mass. The
directed-edge probability measure is
\[
 \mu(u,v)=\frac{w_uw_vc_{uv}}{6T_H(w)}.
\]
Conditional on $u$, its second coordinate has distribution
$\pi_u(v)=w_vc_{uv}/h_u$, and $K_u$ is reversible for $\pi_u$.
We use the constants from that lemma:
\[
 A_C=10+42C^2+2C,
 \qquad B_C=2\sqrt{A_C+1}.
\]
In particular, $K_u(v,y)\le xw_y$, the averaged row entropy is at
most $A_C\log x$, and the average total variation distance between
the rows indexed by the other two vertices of a weighted ordered
triangle is at most $B_C\sqrt{\log x/x}$.

For each non-isolated $u$, apply Lemma~\ref{lem:shared-rejection} to all rows of
$K_u$. Denote the label assigned to $v\in N_H(u)$ by $Y_{u,v}$.
These finite joint laws may be sampled independently for different
vertices $u$.  We will group the neighbors of $u$ according to their
labels, while keeping the total weight of each group at most $e^{\sqrt x}$.
For this purpose, call a label $y$ at the incidence $(u,v)$
\emph{admissible} if
\begin{equation}\label{eq:admissible-label}
 K_u(y,v)\ge w_v e^{-\sqrt x}.
\end{equation}
For a fixed outcome of the labels, define
\[
 A_{u,y}=\{v\in N_H(u):Y_{u,v}=y
       \text{ and }K_u(y,v)\ge w_ve^{-\sqrt x}\}.
\]
The reversed transition in \eqref{eq:admissible-label} is chosen so that
all vertices in one group can be compared with a single probability row:
\begin{equation}\label{eq:label-group-weight}
 w(A_{u,y})
 \le e^{\sqrt x}\sum_{v\in A_{u,y}}K_u(y,v)
 \le e^{\sqrt x}.
\end{equation}
This holds for every outcome of the labels.  The next lemma shows that
the admissibility requirement rejects only a small fraction of incidences
under the triangle distribution.

\begin{lemma}\label{lem:reverse-admissibility}
The probability that $Y_{u,v}$ is inadmissible, averaged over
$(u,v)\sim\mu$, is at most
\begin{equation}\label{eq:inadmissibility-bound}
 \delta_x:=\frac{(A_C+1)\log x}{\sqrt x+\log x}.
\end{equation}
\end{lemma}

\begin{proof}
Under the joint law $(u,v)\sim\mu$, $y\sim K_u(v,\cdot)$,
reversibility gives
\[
 (u,v,y)\ \stackrel{\mathrm{law}}=\ (u,y,v).
\]
Both $K_u(v,y)$ and $K_u(y,v)$ are positive on this law's support:
all $\pi_u$ are positive and
$\pi_u(v)K_u(v,y)=\pi_u(y)K_u(y,v)$.
Consequently
\[
 \ell:=\log\frac{w_v}{K_u(y,v)}
 \quad\text{satisfies}\quad
 \mathbb E\ell=H_j\le A_C\log x.
\]
The density bound applied to the reversed row gives $\ell\ge-\log x$.
Inadmissibility is precisely the event $\ell>\sqrt x$. Applying Markov's
inequality to the nonnegative random variable $\ell+\log x$ proves
\eqref{eq:inadmissibility-bound}. No bound on an individual row's
entropy is needed here.
\end{proof}

We now turn the groups into a graph, for each fixed outcome of the labels.
Make one vertex $(u,y)$ of weight $w_u$ for each nonempty group.
For each edge $uv$ of $H$ whose two incidence labels are admissible,
put an edge between $(u,Y_{u,v})$ and $(v,Y_{v,u})$.
Call the resulting graph $H'$. Figure~\ref{fig:label-splitting} shows
the construction at one vertex.
We take the maximum neighborhood weight of the empty graph to be zero.
\begin{figure}[htbp]
  \centering
  \begingroup
  \definecolor{splitblue}{RGB}{34,92,140}
  \definecolor{splitochre}{RGB}{166,104,24}
  \begin{tikzpicture}[
    x=1cm,y=1cm,
    font=\small,
    vertex/.style={circle,draw=black,fill=white,minimum size=5.5pt,
      inner sep=0pt,line width=.6pt},
    copy/.style={vertex,minimum size=7pt},
    blueedge/.style={draw=splitblue,line width=.9pt},
    ochreedge/.style={draw=splitochre,line width=.9pt},
    every label/.style={font=\small,inner sep=2pt}
  ]
    \node at (1.15,2.22) {Before splitting};
    \node at (7.35,2.22) {After splitting};

    \draw[blueedge] (0,0) -- (2.55,1.55);
    \draw[blueedge] (0,0) -- (2.55,.55);
    \draw[ochreedge] (0,0) -- (2.55,-.55);
    \draw[ochreedge] (0,0) -- (2.55,-1.55);
    \draw[line width=.6pt] (2.55,.55) -- (2.55,1.55);
    \node[vertex] (old-u) at (0,0) {};
    \node[anchor=east,align=center,inner sep=0pt] at (-.17,0)
      {$u$\\[-1pt]$w_u$};
    \node[vertex,label=right:$v_1$] at (2.55,1.55) {};
    \node[vertex,label=right:$v_2$] at (2.55,.55) {};
    \node[vertex,label=right:$v_3$] at (2.55,-.55) {};
    \node[vertex,label=right:$v_4$] at (2.55,-1.55) {};

    \draw[-{Latex[length=2.2mm]},line width=.65pt]
      (3.55,0) -- (4.8,0);
    \node[align=center,font=\footnotesize] at (4.17,.48)
      {split by\\label};

    \draw[blueedge] (6.15,1.05) -- (8.7,1.55);
    \draw[blueedge] (6.15,1.05) -- (8.7,.55);
    \draw[ochreedge] (6.15,-1.05) -- (8.7,-.55);
    \draw[ochreedge] (6.15,-1.05) -- (8.7,-1.55);
    \draw[line width=.6pt] (8.7,.55) -- (8.7,1.55);
    \node[copy,draw=splitblue] at (6.15,1.05) {};
    \node[anchor=east,align=center,inner sep=0pt] at (5.96,1.05)
      {$(u,y)$\\[-1pt]$w_u$};
    \node[copy,draw=splitochre] at (6.15,-1.05) {};
    \node[anchor=east,align=center,inner sep=0pt] at (5.96,-1.05)
      {$(u,z)$\\[-1pt]$w_u$};
    \node[vertex,label=right:$v'_1$] at (8.7,1.55) {};
    \node[vertex,label=right:$v'_2$] at (8.7,.55) {};
    \node[vertex,label=right:$v'_3$] at (8.7,-.55) {};
    \node[vertex,label=right:$v'_4$] at (8.7,-1.55) {};

    \draw[blueedge] (.45,-2.12) -- (.9,-2.12);
    \node[anchor=west,font=\footnotesize] at (.99,-2.12)
      {label $y$ at $u$};
    \draw[ochreedge] (5.75,-2.12) -- (6.2,-2.12);
    \node[anchor=west,font=\footnotesize] at (6.29,-2.12)
      {label $z$ at $u$};
  \end{tikzpicture}
  \endgroup
  \caption{Local view of splitting by admissible labels. The incidences
    $uv_1,uv_2$ receive label $y$, and $uv_3,uv_4$ receive label $z$.
    Each copy of $u$ has weight $w_u$. The remote endpoints are also split:
    $v'_i$ denotes the copy selected by the incidence $v_i u$.
    The displayed triangle survives because its two incidences agree at
    each of its three corners, with all labels admissible. Each retained
    edge is placed once, between its selected copies.}
  \label{fig:label-splitting}
\end{figure}


\begin{lemma}\label{lem:weighted-splitting}
There is a choice of these labels for which the projection
$(u,y)\mapsto u$ preserves weights, maps directed edges injectively
to directed edges of $H$, and satisfies
\begin{align}
 \max_{a\in V(H')}w(N_{H'}(a))&\le e^{\sqrt x},
 \label{eq:split-neighborhood}\\
 T_{H'}(w)&\ge(1-a_C\log x/\sqrt x)T_H(w),
 \label{eq:split-triangles}
\end{align}
where $a_C=6(B_C+A_C+1)$.
\end{lemma}

\begin{proof}
Every old edge creates at most one new edge. The projection therefore
maps directed edges injectively, and the new graph has neither loops
nor multiple edges. At a fixed vertex $(u,y)$, different neighbors
project to different vertices of $H$, because $H$ is simple and each
edge $uv$ was used at most once.  By \eqref{eq:label-group-weight},
\[
 w(N_{H'}(u,y))
 \le w(A_{u,y})
 \le e^{\sqrt x}.
\]
In particular, duplicating a vertex does not duplicate any of its
incident edges at one fixed copy.

An old triangle $\{u,v,z\}$ survives whenever its two labels agree
at each corner and all six incidence labels are admissible. These
conditions are also necessary for its three edges to form a triangle
in $H'$. A triangle of $H'$ projects to three distinct old vertices,
and hence to an old triangle. Moreover, there is at most one triangle
over each old triangle, since each old edge has at most one preimage.
The weights of a surviving triangle are unchanged.

Sample an ordered triangle $(u,v,z)$ with probability
$w_uw_vw_z/(6T_H(w))$. Its ordered-edge marginal is $\mu$.
At the corner $u$, the shared coupling and the row-distance estimate
give
\[
 \mathbb E\,\mathbb P(Y_{u,v}\ne Y_{u,z})
 \le 2B_C\sqrt{\frac{\log x}{x}}.
\]
Each of the other two corners has the same bound, by the permutation
symmetry of the ordered-triangle law. Each of its six directed
incidences has marginal $\mu$, so
Lemma~\ref{lem:reverse-admissibility} bounds the average
inadmissibility probability of each by $\delta_x$.
A union bound, with no independence assumption between these events,
therefore gives
\[
 \frac{\mathbb E(T_H(w)-T_{H'}(w))}{T_H(w)}
 \le6B_C\sqrt{\frac{\log x}{x}}+6\delta_x
 \le a_C\frac{\log x}{\sqrt x}.
\]
The last inequality uses $\log x\ge1$. All label vectors together
have finite support. At least one of them has triangle mass at least
its expectation, proving the claim.
\end{proof}

\subsection{Reduction and iteration}

The preceding construction turns the walk estimates into a reduction
of weighted neighborhoods. Choose $X=X(C)\ge32$ so large that
\begin{equation}\label{eq:iteration-threshold}
 2a_C\frac{\log X}{\sqrt X}\le\frac12.
\end{equation}
Such a choice depends only on $C$.

\begin{proposition}\label{prop:neighborhood-reduction}
Let $G$ be a finite positive-vertex-weighted graph satisfying the
cross-mass condition~\eqref{eq:cross-mass} with constant $C\ge1$, and suppose that
$\max_vw(N_G(v))\le e^x$ for some $x>X(C)$. There is a finite
positive-vertex-weighted graph $G'$ with a weight-preserving map to
$G$ that is injective on directed edges, such that
\begin{align}
 \max_vw(N_{G'}(v))&\le e^{\sqrt x},\label{eq:step-neighborhood}\\
 T_{G'}(w)&\ge
 \left(1-a_C\frac{\log x}{\sqrt x}\right)T_G(w)
       -\frac{M_G(w)}x.\label{eq:step-triangles}
\end{align}
It also satisfies the same cross-mass condition, and $M_{G'}(w)\le M_G(w)$.
\end{proposition}

\begin{proof}
Peel edges whose current common-neighbor weight is less than $1/x$.
If $H$ is the remaining graph, Lemma~\ref{lem:triangle-peeling} gives
\[
 T_H(w)\ge T_G(w)-M_G(w)/x.
\]
Edge deletion does not increase any weighted neighborhood, and it
preserves the cross-mass condition. If $H$ has no edges, take $G'$ to
be empty. Then $T_G(w)\le M_G(w)/x$, which implies
\eqref{eq:step-triangles}.
Otherwise, apply Lemma~\ref{lem:weighted-splitting} to $H$.
Since $x\mapsto\log x/\sqrt x$ is decreasing for $x\ge e^2$,
\eqref{eq:iteration-threshold} ensures
$0\le1-a_C\log x/\sqrt x\le1$. Hence
\begin{align*}
 T_{G'}(w)
 &\ge\left(1-a_C\frac{\log x}{\sqrt x}\right)
           \left(T_G(w)-\frac{M_G(w)}x\right)\\
 &\ge\left(1-a_C\frac{\log x}{\sqrt x}\right)T_G(w)
           -\frac{M_G(w)}x.
\end{align*}
The projection through $H$ is still injective on directed edges.
The edge-mass bound and inherited cross-mass condition follow from
Lemma~\ref{lem:invariant-preservation}.
\end{proof}

\begin{proof}[Proof of Theorem~\ref{thm:weighted-triangles}]
If $G$ has no edges, both sides vanish. Otherwise choose $x_0\ge1$
such that every neighborhood has weight at most $e^{x_0}$. While
$x_i>X$, apply Proposition~\ref{prop:neighborhood-reduction} and
set $x_{i+1}=\sqrt{x_i}$. Write $G_i$ for the successive graphs,
$T_i=T_{G_i}(w)$ and $M_i=M_{G_i}(w)$; the weights at each stage
are the inherited weights of its vertex copies. Stop after $k$ steps,
when $x_k\le X$. The procedure is finite, and the $k=0$ case is
allowed.

We record bounds independent of $x_0$ and $k$. If $k>0$, put
$y=x_{k-1}>X$. The used parameters, in reverse order, are
$y,y^2,y^4,\ldots,y^{2^{k-1}}$. For
$f(t)=\log t/\sqrt t$ and $t\ge X\ge32$,
\[
 \frac{f(t^2)}{f(t)}=\frac2{\sqrt t}\le\frac12.
\]
As $f$ is decreasing on $[X,\infty)$, it follows that
\begin{equation}\label{eq:summable-losses}
 \sum_{i<k}a_Cf(x_i)\le2a_Cf(X)\le\frac12,
 \qquad
 \sum_{i<k}\frac1{x_i}\le\frac{1/X}{1-1/X}\le1.
\end{equation}
Both inequalities also hold for empty sums when $k=0$.

Let $\varepsilon_i=a_Cf(x_i)$. Iterating
\eqref{eq:step-triangles}, using $M_i\le M_0$ and
$0\le1-\varepsilon_i\le1$, gives
\[
 T_k\ge\left(\prod_{i<k}(1-\varepsilon_i)\right)T_0
          -M_0\sum_{i<k}\frac1{x_i}
 \ge\frac12T_0-M_0.
\]
Here $\prod(1-\varepsilon_i)\ge1-\sum\varepsilon_i$.
Every neighborhood of $G_k$ has weight at most $e^X$, so
\[
 3T_k
 =\sum_{\{u,v\}\in E(G_k)}w_uw_vw(N_{G_k}(u)\cap N_{G_k}(v))
 \le e^X M_k\le e^X M_0.
\]
Combining the last two displays proves the theorem with
$C_\triangle(C)=2(1+e^{X(C)}/3)$. All
thresholds and constants depend only on $C$; neither the number of
vertices nor the initial neighborhood bound enters them.
\end{proof}


Applying Theorem~\ref{thm:weighted-triangles} to the cross-mass bound
of Lemma~\ref{lem:cross-mass} gives the variational estimate needed
for the independent-set argument: for every maximizer $w$ on a finite
$K_r$-free graph,
\begin{equation}\label{eq:optimizer-triangle}
 T_G(w)\le B_r M_G(w),\qquad B_r=C_{\triangle}(16r^2).
\end{equation}
The zero-vertex case has both masses zero.
